\section{Invariants}

% ── slides p.2–p.6 ──────────────────────────────────────────
An \textbf{invariant} of a dynamic system is an assertion that holds at every reachable state. Invariants are negative oracles: a state violating the invariant is certainly unreachable, with no enumeration of the reachable set. This chapter develops two structural classes for Petri nets: \emph{place invariants} (S-invariants) and \emph{transition invariants} (T-invariants), computable from the incidence matrix alone.\footnote{Desel J., Esparza J., \emph{Free Choice Petri Nets}, Cambridge University Press, \url{https://www7.in.tum.de/~esparza/bookfc.html}.}

\begin{examplebox}{The mutilated chessboard}
Remove two opposite corners from a chessboard (62 cells remain) and take 31 dominoes, each covering two adjacent cells. The arithmetic fits ($31 \times 2 = 62$), yet \emph{no tiling exists}: the removed corners share a colour, leaving 32 cells of one colour and 30 of the other, while every domino covers one of each: so any tiling covers equal numbers. The conserved quantity (covered-black minus covered-white) discriminates reachable from unreachable configurations without trying a single tiling.
\end{examplebox}

% ── slides p.7–p.19 ──────────────────────────────────────────
An invariant is always relative to a set of admissible moves. For a polygon under \{rotate, move, mirror\} the perimeter and area are invariant.

Adding \emph{scale} kills both, but spares the combinatorial quantities (number of vertices and sides, vertex degrees, convexity) and the label (colour).

Adding \emph{stretch} (non-uniform scaling) still preserves colour, convexity, sides and vertices, since affine maps send edges to edges; degenerate stretches, however, can merge vertices and break degrees. Choosing the move set is part of stating the invariant.

% ── slides p.20–p.33 ──────────────────────────────────────────
For a Petri net the moves are the firings of enabled transitions. Trivial invariants: the labels and the static skeleton $(P, T, F)$ (firing changes only the marking), the token count of places no reachable firing touches. The token count of an arbitrary place is \emph{not} invariant: any transition consuming from it changes it. More interestingly, the system-level properties themselves are invariants:

\begin{notebox}{Lemma: properties closed under reachability}
If $(P, T, F, M_0)$ is live (resp.\ deadlock-free, bounded, cyclic) and $M \in [M_0\rangle$, then $(P, T, F, M)$ is live (resp.\ deadlock-free, bounded, cyclic). Here \emph{cyclic} means $\forall M \in [M_0\rangle.\ M_0 \in [M\rangle$: no point of no return.
\par\smallskip
\textit{Proof shape (same for all four).} Any $M' \in [M\rangle$ is also in $[M_0\rangle$ by transitivity of reachability, since $[M\rangle \subseteq [M_0\rangle$. For live and deadlock-free (purely universal properties) the claim instantiates to every such $M'$ directly. For boundedness, which is $\exists k\, \forall M'.\ M'(p) \leq k$, the same bound $k$ that witnesses $[M_0\rangle$ also bounds the smaller set $[M\rangle$. For cyclicity, concatenate the return path $M' \to^\ast M_0$ with $M_0 \to^\ast M$. \hfill$\square$
\end{notebox}

% ── slides p.34–p.44 ──────────────────────────────────────────
\subsection{S-invariants}

Structural invariants are vectors of rationals computed from the net's skeleton, independent of the initial marking; a basis is computable in polynomial time. Reading a model through its invariants exposes the conservation laws that govern every reachable marking.

\begin{definitionbox}{S-invariant and its fundamental property}
An \textbf{S-invariant} (place invariant) of $N = (P, T, F)$ with incidence matrix $\mathbf{N}$ is a row vector $\mathbf{I}$ of length $|P|$ with
\[
  \mathbf{I} \cdot \mathbf{N} = \mathbf{0}.
\]
$\mathbf{I}$ assigns a weight to each place; the equation says every transition preserves the weighted token count. Linear combinations of S-invariants are S-invariants (the solution set of a homogeneous system is a vector space, so a basis describes them all).
\par\smallskip
\textbf{Fundamental property.} For every S-invariant $\mathbf{I}$ and every $M \in [M_0\rangle$: $\;\mathbf{I} \cdot M = \mathbf{I} \cdot M_0$.
\par\smallskip
\textit{Proof.} $M = M_0 + \mathbf{N} \cdot \vec{\sigma}$ for some firing sequence $\sigma$ (marking equation), so $\mathbf{I} \cdot M = \mathbf{I} \cdot M_0 + (\mathbf{I} \cdot \mathbf{N}) \cdot \vec{\sigma} = \mathbf{I} \cdot M_0$. \hfill$\square$
\end{definitionbox}

On the vending-machine net of the incidence-matrices chapter ($M_0 = 4p_1 + p_3$), the system $\mathbf{I} \cdot \mathbf{N} = \mathbf{0}$ expands to $x_1 = x_2$, $x_3 = x_4 = x_5$: every S-invariant has the form $[n, n, m, m, m]$, a two-dimensional space with basis $\{[1,1,0,0,0],\, [0,0,1,1,1]\}$: the four stock-side tokens stay in $\{p_1, p_2\}$, the single control token stays in $\{p_3, p_4, p_5\}$.

\begin{notebox}{A cheap necessary condition for reachability}
Reachability is decidable but \textsc{ExpSpace}-hard. Two cheap necessary tests: if some S-invariant has $\mathbf{I} \cdot M \neq \mathbf{I} \cdot M_0$, then $M \notin [M_0\rangle$; and if $\mathbf{N} \cdot \mathbf{x} = M - M_0$ has no rational solution, then $M \notin [M_0\rangle$. Neither is sufficient. Example: in the producer/consumer net, the consumer-loop invariant $\mathbf{I} = \textsf{cons\_free} + \textsf{cons\_busy}$ has a fixed value determined by $M_0$; any candidate marking with a different consumer count (say one consumer-state token where $M_0$ had none) is rejected as unreachable without exploring a single state.
\end{notebox}

% ── slides p.46–p.66 ──────────────────────────────────────────
The conservation reading: tokens are coins, places are denominations, $\mathbf{I}(p)$ is the value of denomination $p$; every firing is a change-making operation that preserves the total value $\mathbf{I} \cdot M$. The same intuition (transitions redistribute but never create or destroy the weighted substance) yields a local test:

\begin{notebox}{Alternative characterisation}
$\mathbf{I} : P \to \mathbb{Q}$ is an S-invariant iff for every transition $t$:
\[
  \sum_{p \in {}^\bullet t} \mathbf{I}(p) \;=\; \sum_{p \in t^\bullet} \mathbf{I}(p)
  \qquad \text{(flow-in weight = flow-out weight)}.
\]
\textit{Proof.} $\mathbf{I} \cdot \mathbf{N} = \mathbf{0}$ holds iff $\mathbf{I} \cdot \mathbf{t} = 0$ for every column $\mathbf{t}$; expanding with $\mathbf{N}(p,t) = +1$ on $t^\bullet$, $-1$ on ${}^\bullet t$, $0$ elsewhere gives $\sum_{p \in t^\bullet} \mathbf{I}(p) - \sum_{p \in {}^\bullet t} \mathbf{I}(p) = 0$. \hfill$\square$
\par\smallskip
One balance per transition, checkable on the picture without the matrix: the workhorse both for verifying and for \emph{constructing} invariants.
\end{notebox}

% ── slides p.67–p.76 ──────────────────────────────────────────
\begin{examplebox}{Bike rental: which candidates are S-invariants?}
\begin{center}
\begin{tikzpicture}[node distance=0.5cm and 0.8cm]
  \node[bpmn event, label=above:{\scriptsize $p_1$ persons}] (p1) {};
  \fill[accent] (p1.center) circle (1.6pt);
  \node[bpmn event, label=below:{\scriptsize $p_2$ bikes}] (p2) at ($(p1)+(0,-1.3)$) {};
  \fill[accent] (p2.center) circle (1.6pt);
  \node[bpmn task, rounded corners=0pt, font=\tiny\sffamily] (take) at ($(p1)!0.5!(p2)+(1.6,0)$) {take};
  \node[bpmn event, label=right:{\scriptsize $p_3$ riders}] (p3) at ($(take)+(1.5,0)$) {};
  \node[bpmn task, rounded corners=0pt, font=\tiny\sffamily] (leave) at ($(p1)!0.5!(p2)+(1.6,1.4)$) {leave};
  \draw[bpmn flow] (p1) -- (take); \draw[bpmn flow] (p2) -- (take);
  \draw[bpmn flow] (take) -- (p3);
  \draw[bpmn flow] (p3) to[bend right=25] (leave);
  \draw[bpmn flow] (leave) to[bend right=20] (p1);
  \draw[bpmn flow] (leave) to[bend left=35] (p2.west);
\end{tikzpicture}
\end{center}
With \textsf{take}$: p_1 + p_2 \to p_3$ and \textsf{leave}$: p_3 \to p_1 + p_2$, the local balance reduces to one equation: $\mathbf{I}(p_1) + \mathbf{I}(p_2) = \mathbf{I}(p_3)$. Testing candidates $[p, b, r]$: $[1,0,1]$ \textbf{yes} (persons free or riding: constant), $[0,1,1]$ \textbf{yes} (bikes free or ridden), $[1,-1,0]$ \textbf{yes} (persons-minus-bikes surplus), $[1,1,2]$ \textbf{yes} (a rider counts as person + bike); $[1,1,-1]$, $[1,1,1]$, $[1,2,2]$ \textbf{no} (balance fails; $[1,1,1]$ double-counts the rider). The four winners are not independent: $[1,1,2] = [1,0,1] + [0,1,1]$ and $[1,-1,0] = [1,0,1] - [0,1,1]$: a basis is $\{[1,0,1], [0,1,1]\}$, the two natural conservation laws.
\end{examplebox}

% ── slides p.77–p.101 ──────────────────────────────────────────
\begin{examplebox}{Reading invariants off the picture}
\emph{Traffic lights with mutex}: per light, the three state places form a uniform invariant (``each light shows exactly one colour''); a third invariant weights green and yellow of both lights together with the mutex place: the safety law ``at most one light is away from red''. Sums of these are again invariants.
\par\smallskip
\emph{Producer/consumer}: by the linear system, or directly on the picture: weight the producer loop $\{\textsf{prod\_free}, \textsf{prod\_busy}\}$ with 1; the balance at \textsf{prod\_end} then forces the buffer weight to 0; zero-weight places decouple the picture, and the consumer loop independently takes weight 1. Basis: $[1,1,0,0,0]$ and $[0,0,0,1,1]$, the producer and the consumer never disappear, whatever the buffer does. Their sum is the combined law ``one producer-state plus one consumer-state token, always''.
\par\smallskip
\emph{Dining philosophers}: each pair $(\textsf{think}_i, \textsf{eat}_i)$ is a two-place loop, giving the uniform invariant $\textsf{think}_i + \textsf{eat}_i = 1$ (five of them). For the resources, the naive ``$\textsf{eat}_i$ + both adjacent forks'' with unit weights fails the balance at $\textsf{start\_eating}_i$ (flow-in 2, flow-out 1); the corrected versions are either $2\,\textsf{eat}_i + f_{\text{left}} + f_{\text{right}}$, or (cleanest) one invariant \emph{per fork}: $f + \textsf{eat}_A + \textsf{eat}_B = 1$ for the two neighbours $A, B$ of $f$, i.e.\ each fork is free or held by exactly one neighbour. Sums build global laws (``five philosophers, each in some state''); the picture-level method (one local balance per transition, small loops first, weight 0 for outsiders) scales where solving the system would be tedious.
\end{examplebox}

% ── slides p.102–p.129 ──────────────────────────────────────────
\subsection{S-invariants and system properties}

\begin{definitionbox}{Semi-positive, support, positive, uniform}
An S-invariant $\mathbf{I}$ is \textbf{semi-positive} ($\mathbf{I} > \mathbf{0}$) when $\mathbf{I} \geq \mathbf{0}$ and $\mathbf{I} \neq \mathbf{0}$; its \textbf{support} is $\langle \mathbf{I} \rangle = \{p \mid \mathbf{I}(p) > 0\}$; it is \textbf{positive} ($\mathbf{I} \succ \mathbf{0}$) when $\langle \mathbf{I} \rangle = P$; \textbf{uniform} when semi-positive with coefficients in $\{0, 1\}$ (a plain token count). Every semi-positive S-invariant satisfies ${}^\bullet\langle \mathbf{I} \rangle = \langle \mathbf{I} \rangle^\bullet$: the transitions feeding the support are exactly those draining it (symmetrically for T-invariants).
\end{definitionbox}

\begin{notebox}{Theorem: positive S-invariant $\Rightarrow$ bounded}
If the system admits a positive S-invariant $\mathbf{I}$, it is bounded; more precisely every $p \in \langle \mathbf{I} \rangle$ of a semi-positive invariant satisfies
\[
  M(p) \;\leq\; \frac{\mathbf{I} \cdot M_0}{\mathbf{I}(p)} \qquad \text{for every } M \in [M_0\rangle.
\]
\textit{Proof.} $\mathbf{I}(p) M(p) \leq \mathbf{I} \cdot M = \mathbf{I} \cdot M_0$ since the other terms are non-negative; divide by $\mathbf{I}(p) > 0$. \hfill$\square$
\end{notebox}

On the bike rental with one person and one bike, $\mathbf{I} = [1,1,2]$ is positive with $\mathbf{I} \cdot M_0 = 2$, so $M(p_3) \leq 2/2 = 1$: at most one rider, ever. The sharper $\mathbf{I} = [1,0,1]$ gives the same bound from $1/1$: different invariants give different bounds, and hunting the tightest is a linear-programming question. On the vending machine, $[1,1,1,1,1]$ (sum of the basis) is positive with $\mathbf{I} \cdot M_0 = 5$: every place is 5-bounded.

\begin{notebox}{Theorem: a necessary condition for liveness}
If the system is live, then every semi-positive S-invariant satisfies $\mathbf{I} \cdot M_0 > 0$.
\par\smallskip
\textit{Proof.} Pick $p \in \langle \mathbf{I} \rangle$ and a transition $t$ adjacent to $p$. By liveness some reachable $M$ fires $t$: if $t$ produces into $p$ the successor marks $p$; if $t$ consumes from $p$ then $M$ itself marks $p$. Either way some reachable $M^*$ has $\mathbf{I} \cdot M^* \geq \mathbf{I}(p) M^*(p) > 0$, and the fundamental property transports the value back to $M_0$. \hfill$\square$
\par\smallskip
\textbf{Contrapositive certificate}: a semi-positive S-invariant with $\mathbf{I} \cdot M_0 = 0$ proves the system non-live. To find one, look for an invariant whose support avoids all initially marked places: on the bike rental started with bikes only, $[1,0,1]$ has value 0: \textsf{take} can never fire; on a seven-place exercise net marked only on $p_5, p_6$, the invariant supported on $\{p_2, p_3, p_4\}$ does the job; on the producer/consumer started without consumer tokens, the consumer-loop invariant $[0,0,0,1,1]$ kills liveness.
\end{notebox}

Two markings \emph{agree on all S-invariants} when $\mathbf{I} \cdot M = \mathbf{I} \cdot M'$ for every $\mathbf{I}$: equivalently, when $\mathbf{N} \cdot \mathbf{y} = M' - M$ has a rational solution. Agreement is necessary for mutual reachability, not sufficient: unreachable look-alikes exist. S-invariants are negative oracles: positive invariant $\Rightarrow$ bounded; $\mathbf{I} \cdot M_0 = 0$ (semi-positive) $\Rightarrow$ non-live; $\mathbf{I} \cdot M \neq \mathbf{I} \cdot M_0$ $\Rightarrow$ $M$ unreachable, and none of the converses holds. The producer/consumer with unbounded buffer shows the first gap concretely: the balance at \textsf{prod\_end} forces buffer weight $x = 0$ in every invariant ($y = y + x$), so no positive S-invariant exists: consistently with the unbounded buffer.

% ── slides p.130–p.145 ──────────────────────────────────────────
\subsection{T-invariants}

\begin{definitionbox}{T-invariant and its fundamental property}
A \textbf{T-invariant} of $N$ is a column vector $\mathbf{J}$ of length $|T|$ with
\[
  \mathbf{N} \cdot \mathbf{J} = \mathbf{0}.
\]
$\mathbf{J}(t)$ is an occurrence count for $t$: firing every transition $\mathbf{J}(t)$ times causes zero net change everywhere.
\par\smallskip
\textbf{Fundamental property.} For $M \xrightarrow{\sigma} M'$: $\vec{\sigma}$ is a T-invariant iff $M = M'$. \textit{Proof}: immediate from the marking equation $M' = M + \mathbf{N} \cdot \vec{\sigma}$, both directions. \hfill$\square$
\par\smallskip
\textbf{Alternative characterisation} (dual of the S case): $\mathbf{J}$ is a T-invariant iff at every \emph{place} $p$, $\sum_{t \in {}^\bullet p} \mathbf{J}(t) = \sum_{t \in p^\bullet} \mathbf{J}(t)$: weighted production equals weighted consumption, one balance per place, readable off the picture.
\end{definitionbox}

% Corrected 2026-06-12: the candidate-testing exercise was wrongly attributed
% to the vending machine (whose T-invariants are [a,a,a+b,a,b], so none of the
% listed candidates qualifies) and [1,1,2,1,2] was marked as an invariant. The
% slides (p.137-145) pose it on a different five-place exercise net, restated
% here by its arcs; verdicts recomputed: yes [1,0,0,1,1], [1,1,1,1,2],
% [0,1,1,0,1]; no [1,1,2,1,2] (p3 balance 3 vs 2), no [1,1,1,0,2] (p2: 1 vs 0).
T-invariants are the reproducibility certificates: multisets of firings that, when executable, return the marking to its starting point: cycles of the reachability graph up to permutation. On the bike rental, $[1,1]$ (one \textsf{take}, one \textsf{leave}) is a T-invariant. A richer exercise net: $t_1 : p_1 \to p_2 + p_3$, $t_2 : p_1 \to p_4 + p_5$, $t_4 : p_2 \to p_4$, $t_3 : p_5 \to p_3$, $t_5 : p_3 + p_4 \to p_1$: a token leaves $p_1$ along one of two routes, splits, and $t_5$ reassembles it. Testing candidates $[t_1, \ldots, t_5]$ against the per-place balance: $[1,0,0,1,1]$, $[1,1,1,1,2]$ and $[0,1,1,0,1]$ are T-invariants; $[1,1,2,1,2]$ is not (at $p_3$, in-flow $\mathbf{J}(t_1) + \mathbf{J}(t_3) = 3$ against out-flow $\mathbf{J}(t_5) = 2$), nor is $[1,1,1,0,2]$ (the zero on $t_4$ breaks the balance at $p_2$).

% ── slides p.146–p.158 ──────────────────────────────────────────
\subsection{T-invariants and system properties}

\begin{notebox}{Reproduction lemma and the bounded-live theorem}
\emph{Pigeonhole}: a path through more states than exist must revisit one.
\par\smallskip
\textbf{Lemma.} In a bounded system, if $M_0 \xrightarrow{\sigma}$ for an infinite $\sigma$, there is a semi-positive T-invariant supported on transitions of $\sigma$.
\textit{Proof.} The reachable set is finite, so the infinite trajectory revisits a marking: $M_i = M_j$ with $i < j$; the Parikh vector of the segment $t_{i+1} \cdots t_j$ is a non-zero T-invariant by the fundamental property. \hfill$\square$
\par\smallskip
\textbf{Theorem.} A bounded \emph{and} live system has a \emph{positive} T-invariant.
\textit{Proof.} By liveness, from any marking a sequence firing \emph{every} transition at least once exists; chain $k + 1$ such sequences ($k = |[M_0\rangle|$): among the $k + 2$ intermediate markings two coincide, and the enclosed segment is a T-invariant containing at least one full sweep: positive on every transition. \hfill$\square$
\end{notebox}

The contrapositive chain is the diagnostic: no positive T-invariant $\Rightarrow$ not both live and bounded; combined with liveness it proves unboundedness, with boundedness (e.g.\ via a positive S-invariant) it proves non-liveness. The implication does \emph{not} reverse: a positive T-invariant alone certifies nothing:
\begin{itemize}
  \item two dead transitions over two empty places admit $\mathbf{J} = [1,1]$, yet nothing ever fires;
  \item the producer/consumer with an unmarked consumer side has $\mathbf{J} = [1,1,1,1]$ positive, yet is non-live (dead consumer) \emph{and} unbounded (the producer pumps the buffer);
  \item variants where the buffer is never consumed are live but unbounded, with the same positive $\mathbf{J}$.
\end{itemize}

% ── slides p.159–p.163 ──────────────────────────────────────────
\subsection{Strong connectedness via invariants}

\begin{notebox}{Theorem (proof omitted)}
A weakly connected net with a positive S-invariant \emph{and} a positive T-invariant is strongly connected. Contrapositive: a weakly-but-not-strongly connected net cannot have both.
\end{notebox}

The producer/consumer illustrates one direction: not strongly connected, positive T-invariant $[1,1,1,1]$, but no positive S-invariant (the buffer is stuck at weight 0). A mirror example with a positive S-invariant (weights $\geq 1$ everywhere, place $p_1$ at 2) but no positive T-invariant is bounded yet non-live: two invariant computations pin down the qualitative behaviour with no simulation at all.
